\documentclass[journal,twoside]{IEEEtran}
\usepackage{cite}
\usepackage{amsmath,amssymb,amsfonts}
\usepackage{algorithmic}
\usepackage{graphicx}
\usepackage{textcomp}
\usepackage{xcolor}
\usepackage{booktabs}
\usepackage{hyperref}

\begin{document}

\title{An Adaptive Hybrid Quantum-Resistant Cloud Storage Architecture with Zero-Trust Key Splitting and Cryptographic Hash-Chained Forensics}

\author{Author Name,~\IEEEmembership{Member,~IEEE,} and Co-Author Name,~\IEEEmembership{Senior Member,~IEEE}%
\thanks{This research was supported by Cloud Security and Quantum-Resistant Infrastructure Research Initiatives.}}

\markboth{IEEE Transactions on Cloud Computing,~Vol.~XX, No.~X, August~2026}%
{Author \MakeLowercase{\textit{et al.}}: Adaptive Hybrid Quantum-Resistant Cloud Storage Architecture}

\maketitle

\begin{abstract}
The rapid emergence of quantum computing presents an existential threat to modern cloud storage security. Classical asymmetric cryptosystems—including RSA and Elliptic Curve Cryptography (ECC)—rely on prime factorization and discrete logarithms, which will be solved in polynomial time by quantum computers running Shor's Algorithm. Consequently, adversaries are actively engaged in "Harvest Now, Decrypt Later" (HNDL) attacks, intercepting encrypted cloud data today to retroactively decrypt it once quantum hardware matures. To address this threat without incurring unsustainable cloud performance penalties, this paper proposes an \textbf{Adaptive Hybrid Quantum-Resistant Cloud Storage System}. The architecture introduces five core technological innovations: (1) an \textbf{Adaptive Encryption Decision Engine} that dynamically selects between ECC-256, CRYSTALS-Kyber-512, or a Hybrid mode based on data sensitivity and payload size; (2) a \textbf{Dual-Key Hybrid KEM Scheme} pairing ECDH (SECP256R1) with CRYSTALS-Kyber-512 via HKDF key derivation to ensure dual-algorithm fail-safe security; (3) a \textbf{Zero-Trust Shamir Secret Sharing (2-of-3)} scheme over $\text{GF}(257)$ that fragments master keys across isolated storage domains (Relational DB + File Package Header) to eliminate single points of key compromise; (4) an \textbf{Early-Drop Outer HMAC-SHA256 Authentication Layer} that rejects tampered ciphertexts prior to AES-256-GCM decryption; and (5) an \textbf{Immutable Cryptographic Hash-Chain Audit Log} integrated with an automated \textbf{Real-Time Anomaly Engine}. Empirical evaluations across 100~KB, 1~MB, and 5~MB payloads confirm that the proposed Hybrid scheme incurs a negligible latency overhead of only ~0.16~ms over classical ECC for a 1~MB file, demonstrating enterprise feasibility for post-quantum cloud environments.
\end{abstract}

\begin{IEEEkeywords}
Post-Quantum Cryptography (PQC), CRYSTALS-Kyber, Elliptic Curve Cryptography (ECC), Hybrid KEM, Shamir Secret Sharing, Adaptive Encryption, Zero-Trust Architecture, Hash-Chain Audit Log, Cloud Storage Security.
\end{IEEEkeywords}

\section{Introduction}
\IEEEPARstart{C}{loud} computing infrastructures store vast quantities of sensitive enterprise, financial, medical, and national defense data. To preserve confidentiality, these platforms depend on public-key cryptosystems such as RSA-3072 and Elliptic Curve Cryptography (ECC-256). However, the eventual deployment of cryptographically relevant quantum computers running Shor's Algorithm will reduce the time complexity of solving prime factorizations and discrete logarithms from sub-exponential to polynomial time $\mathcal{O}(n^3)$ \cite{shor1994}.

This vulnerability has popularized the \textbf{"Harvest Now, Decrypt Later" (HNDL)} attack vector. Adversaries routinely intercept and archive encrypted cloud payloads today, intending to decrypt them retroactively once quantum computers become operational. Because medical histories, defense secrets, and intellectual property require confidentiality spans of 30 to 50+ years, post-quantum migration must occur immediately.

To address this challenge, the National Institute of Standards and Technology (NIST) standardized post-quantum cryptographic primitives, including **CRYSTALS-Kyber** (ML-KEM) \cite{bos2018kyber}. Nevertheless, deploying pure post-quantum algorithms introduces two significant operational hurdles:
\begin{enumerate}
    \item \textbf{Novel Mathematical Risk}: Lattice-based mathematics underlying Kyber is relatively young compared to decades of ECC cryptanalysis. Should an undiscovered mathematical vulnerability be exposed in Kyber, systems relying solely on Kyber would face immediate compromise.
    \item \textbf{Performance and Storage Overhead}: Kyber public keys (800 bytes) and ciphertexts (768 bytes) are significantly larger than ECC-256 (64 bytes). Applying heavy quantum encryption indiscriminately to all routine data degrades cloud throughput and increases storage costs.
\end{enumerate}

To solve these dual challenges, this paper presents a comprehensive, production-ready \textbf{Adaptive Hybrid Quantum-Resistant Cloud Storage System}.

\section{System Architecture \& Methodological Design}
The architecture integrates five cryptographically decoupled components orchestrating file ingress and egress.

\subsection{Adaptive Encryption Decision Engine}
Rather than applying a monolithic encryption standard, the system evaluates payload size ($S$) and user-specified sensitivity ($\sigma \in \{\text{Low}, \text{Medium}, \text{High}\}$):
\begin{equation}
\text{Mode} = \begin{cases} 
\text{ECC-256}, & \text{if } \sigma = \text{Low} \\
\text{Kyber-512}, & \text{if } \sigma = \text{Medium} \lor (S > 5\text{MB} \land \sigma = \text{High}) \\
\text{Hybrid}, & \text{if } \sigma = \text{High} \land S \le 5\text{MB}
\end{cases}
\end{equation}
This policy ensures routine data remains computationally lightweight while high-value assets receive maximum hybrid security.

\subsection{Dual-Key Hybrid KEM and HKDF Expansion}
In Hybrid mode, parallel key exchanges generate classical secret $S_{\text{ecc}}$ via ECDH (SECP256R1) and post-quantum secret $S_{\text{kyber}}$ via Kyber-512 KEM encapsulation:
\begin{equation}
S_{\text{combined}} = S_{\text{ecc}} \mathbin{\Vert} S_{\text{kyber}}
\end{equation}
The joint secret is expanded using HKDF-SHA256 with a 16-byte random salt and context string $b\text{'quantum\_hybrid\_kdf'}$:
\begin{equation}
(K_{\text{enc}} \mathbin{\Vert} K_{\text{auth}}) = \text{HKDF-SHA256}(S_{\text{combined}}, \text{salt}, \text{info}, \text{len}=64)
\end{equation}
where $K_{\text{enc}}$ (32 bytes) serves as the AES-256 payload key and $K_{\text{auth}}$ (32 bytes) serves as the outer authentication key.

\subsection{Zero-Trust Shamir Secret Sharing over $\text{GF}(257)$}
To eliminate single points of key compromise, master key $K_{\text{enc}}$ is split using a $(k=2, n=3)$ threshold polynomial over Galois Field $\text{GF}(257)$ for each byte $S \in [0, 255]$:
\begin{equation}
f(x) = (S + a_1 x) \pmod{257}, \quad a_1 \in_R [1, 256]
\end{equation}
Three shares are evaluated at $x \in \{1, 2, 3\}$. Fragment 1 is encrypted and attached to the file payload header in object storage, while Fragment 2 is encrypted and stored in the SQLite database (`key_fragments` table). Reconstructing $S$ at $x=0$ uses Lagrange Interpolation:
\begin{equation}
f(0) = \sum_{i=1}^{2} y_i \prod_{j \neq i} \frac{-x_j}{x_i - x_j} \pmod{257}
\end{equation}
An adversary breaching only the database obtains Fragment 2, which mathematically yields zero information regarding $K_{\text{enc}}$.

\subsection{Early-Drop Outer HMAC Double Authentication}
The AES-256-GCM encrypted payload is wrapped with an outer HMAC signature:
\begin{equation}
\text{MAC} = \text{HMAC-SHA256}(K_{\text{auth}}, \text{Nonce} \mathbin{\Vert} \text{Ciphertext})
\end{equation}
Upon download request, the outer MAC is verified at the API boundary before key reconstruction or decryption, dropping tampered payloads immediately to prevent CPU exhaustion DoS attacks.

\subsection{Immutable Cryptographic Hash-Chain Log & Anomaly Engine}
All system events are logged to a tamper-evident audit chain:
\begin{equation}
\text{Hash}_N = \text{SHA256}(\text{Hash}_{N-1} \mathbin{\Vert} \text{JSON}(\text{Event Data}))
\end{equation}
Recalculating the hash chain detects historical log deletion or alteration. If $\ge 5$ decryption failures occur for a file/user ID, an automated Anomaly Alert is triggered.

\section{Experimental Evaluation and Results}
The system was evaluated over synthetic file payloads across three size categories ($100\text{ KB}$, $1\text{ MB}$, $5\text{ MB}$) with 20 iterations per benchmark suite.

\begin{table}[htbp]
\caption{Cryptographic Overhead and Storage Metrics}
\label{tab:benchmarks}
\centering
\begin{tabular}{lcccccc}
\toprule
\textbf{Cryptographic Scheme} & \textbf{KeyGen (ms)} & \textbf{Enc (ms)} & \textbf{Dec (ms)} & \textbf{PK (B)} & \textbf{CT (B)} & \textbf{Quantum Safe?} \\
\midrule
RSA-3072 (Baseline) & 4.52 & 0.15 & 3.21 & 384 & 384 & No \\
ECC-256 (Classical) & 0.98 & 0.44 & 0.46 & 64 & 0 & No \\
Kyber-512 (Post-Quantum) & 0.02 & 0.03 & 0.04 & 800 & 768 & Yes \\
Kyber-768 & 0.03 & 0.04 & 0.05 & 1184 & 1088 & Yes \\
Kyber-1024 & 0.04 & 0.05 & 0.06 & 1568 & 1568 & Yes \\
\textbf{Proposed Hybrid} & \textbf{0.31} & \textbf{0.48} & \textbf{0.49} & \textbf{864} & \textbf{768} & \textbf{Yes} \\
\bottomrule
\end{tabular}
\end{table}

\begin{table}[htbp]
\caption{Payload Processing Latency Across File Sizes (20 Iterations)}
\label{tab:sizes}
\centering
\begin{tabular}{lcccc}
\toprule
\textbf{Scheme} & \textbf{100 KB Latency} & \textbf{1 MB Latency} & \textbf{5 MB Latency} & \textbf{PQC Safe?} \\
\midrule
ECC-256 & 7.83 ms & 3.51 ms & 13.19 ms & No \\
Kyber-512 & 0.21 ms & 2.99 ms & 10.94 ms & Yes \\
\textbf{Proposed Hybrid} & \textbf{1.33 ms} & \textbf{3.67 ms} & \textbf{12.19 ms} & \textbf{Yes} \\
\bottomrule
\end{tabular}
\end{table}

\section{Security Threat Analysis}
\begin{table}[htbp]
\caption{Threat Vector Resilience Comparison}
\label{tab:threats}
\centering
\begin{tabular}{lccc}
\toprule
\textbf{Attack Vector} & \textbf{ECC-256} & \textbf{Kyber-512} & \textbf{Proposed Hybrid System} \\
\midrule
Classical Brute Force & Resistant & Resistant & \textbf{Highly Resistant} \\
Shor's Quantum Algorithm & Vulnerable & Resistant & \textbf{Resistant (Quantum Safe)} \\
Novel Kyber Math Flaw & Resistant & Vulnerable & \textbf{Resistant (ECC Fallback)} \\
Database Key Theft & Vulnerable & Vulnerable & \textbf{Resistant (Shamir 2-of-3)} \\
DoS CPU Exhaustion & Vulnerable & Vulnerable & \textbf{Resistant (Outer HMAC)} \\
Log Erasure / Tampering & Vulnerable & Vulnerable & \textbf{Resistant (Hash Chain)} \\
\bottomrule
\end{tabular}
\end{table}

\section{Conclusion}
This paper presented an Adaptive Hybrid Quantum-Resistant Cloud Storage System combining classical ECC with NIST CRYSTALS-Kyber via HKDF-SHA256 key derivation, Shamir Secret Sharing key isolation, outer HMAC double authentication, and hash-chained audit logging. Empirical benchmarks prove that the Hybrid scheme adds only ~0.16~ms overhead for a 1~MB payload while delivering complete quantum immunity, zero-trust key isolation, and tamper-evident auditability.

\begin{thebibliography}{1}
\bibitem{shor1994} P. W. Shor, ``Algorithms for quantum computation: discrete logarithms and factoring,'' in \textit{Proc. 35th Annual Symp. Foundations of Computer Science}, 1994, pp. 124--134.
\bibitem{bos2018kyber} J. Bos \textit{et al.}, ``CRYSTALS-Kyber: a CCA-secure module-lattice-based KEM,'' in \textit{IEEE European Symp. Security and Privacy (EuroS\&P)}, 2018, pp. 253--267.
\bibitem{shamir1979} A. Shamir, ``How to share a secret,'' \textit{Commun. ACM}, vol. 22, no. 11, pp. 612--613, 1979.
\bibitem{hkdf} H. Krawczyk, ``Cryptographic extraction and key derivation: The HKDF scheme,'' in \textit{Proc. CRYPTO}, 2010, pp. 631--648.
\end{thebibliography}

\end{document}
